\section{Bugs found and fixed}
\label{sec:bugs}

This section is append-only: entries stay even after being fixed, as
the record of what was wrong, how it was found, and why the fix is
correct.

\subsection{2026-07-16: electron mass used instead of proton mass in the recombination cost}
\label{sec:mass-bug}

\texttt{radiation\_get\_part\_rate\_to\_fully\_ionize}
(\texttt{src/feedback/GEAR/radiation.c}) computed the electron number
density as
\begin{equation}
  n_e = \frac{X_H \rho}{\mu\, m_e} \qquad (\text{wrong: } m_e = \text{electron mass}),
\end{equation}
instead of Eq.~\ref{eq:recomb-cost}'s $n_e = X_H\rho/m_p$. This
inflated every gas particle's recombination cost -- and hence the
budget a star was charged to maintain it as ionized -- by a factor of
$m_p/m_e \approx 1836$ (times a further $O(1$--$2\times)$ factor from
the spurious $\mu$ described below). At that inflation, a star's
entire photon budget could not fully afford even one gas particle's
cost, so every rebuild cycle tagged at most one particle (via the
deterministic boundary-particle override), then reported the budget
exhausted.

\textbf{How it was found:} the example
(\texttt{StromgrenSphereClumpSmith2021}) was reported to show
essentially no HII region growth. A single-star, no-clump diagnostic
tracking \texttt{IsIonizedFlags} per particle over time showed exactly
one particle tagged per rebuild cycle, a different one each time
(hot-but-untagged and tagged-but-cold particles coexisting at every
snapshot -- flickering, not a growing core). Given the equilibrium
Str\"omgren volume at these parameters should sustain
$\dot N_{\rm ion}/\Delta\dot N_{\rm particle} \approx 20$ particles
simultaneously, a 20$\times$ shortfall pointed at the cost calculation
itself. Tracing \texttt{radiation\_get\_part\_rate\_to\_fully\_ionize}
found the $m_e$-for-$m_p$ swap.

\textbf{Fix:} \texttt{n\_e = (X\_H * rho) / m\_p}. Committed
\texttt{356ec0cda} (first pass, together with removing $\mu$; see
below for why $\mu$ went back and forth).

\textbf{Verification:} rebuilt, reran the same single-star diagnostic.
Tagged-particle count went from pinned at 1 to a smoothly growing
population (17$\to$59 over $t=0$--0.98\,Myr); $r_{\rm HII}$ grew
continuously with no plateau (0$\to$4.9\,pc over the same window,
still accelerating). \texttt{testRadiationRebuildCheck} and
\texttt{test27cellsStars -n 3 -N 3 -r 5} pass;
\texttt{StromgrenSphere} acceptance runs at \texttt{gas\_mass=0.1} and
\texttt{gas\_mass=0.01} both ran clean to completion with no
regression.

\subsection{2026-07-16: mean-molecular-weight factor -- kept, then removed, after checking the source}
\label{sec:mu-bug}

The first fix pass also removed a $\mu$ factor from both $N_H$ and
$n_e$, on the grounds that a hydrogen-atom count has no natural place
for a mean-particle-mass term. Darwin pointed out that
\citet{Hopkins2018}, Appendix E, states explicitly
\[
  \Delta\dot N_b = N(H)_b\,\beta\,n_{e,b}, \qquad
  N(H)_b = \frac{X_H m_b}{\mu\, m_p},
\]
i.e.\ FIRE-2's own formula \emph{does} include $\mu$. The fix was
reverted to keep $\mu$, pending a check.

Fetching the actual FIRE-2 PDF and re-deriving from the underlying
physics resolved it: the recombination rate within a gas element must
equal $\alpha_B\, n_e\, n_H\, V$, so $N_H \equiv n_H V$ is definitionally
the raw hydrogen atom count $X_H m/m_p$ -- there is no room for $\mu$
in that derivation regardless of what FIRE-2's text says. Both
\citet{Hu2017} (``the number of hydrogen \emph{nuclei}'', explicitly
$N_H = m_g X_H/m_p$) and \citet{Smith2021}
($R_{\rm rec}=\beta m\rho(X_H/m_p)^2$) use the $\mu$-free form, and
\texttt{StromgrenSphereClumpSmith2021} exists specifically to match
those two papers' numbers. $\mu$ was removed again (final state,
Eq.~\ref{eq:recomb-cost}), committed as part of \texttt{356ec0cda}.

\textbf{Method note:} this flip-flopped once before landing on the
derivation-based answer. When a paper's prose labels a quantity in a
way that doesn't match its role in the surrounding equation (FIRE-2
calls $N(H)_b$ ``the number of H atoms'' despite the $\mu$ factor
making it something else), trust the derivation over the label, and
check against multiple independent sources before committing.
